\documentclass[11pt,a4paper]{article}

% -------------------------------------------------
% Packages
% -------------------------------------------------
\usepackage[margin=1in]{geometry}
\usepackage{amsmath,amssymb,amsfonts,mathtools}
\usepackage{bm}
\usepackage{physics}
\usepackage{enumitem}
\usepackage{hyperref}

% -------------------------------------------------
% Commands
% -------------------------------------------------
\newcommand{\R}{\mathbb{R}}
\newcommand{\dd}{\mathrm{d}}
\newcommand{\PP}{\mathcal{P}}
\newcommand{\QQ}{\mathcal{Q}}

\begin{document}

\title{2D Hindmarsh--Rose Model}
\author{}
\date{}

\begin{document}

\maketitle

\section{The Simplified Model}

The general 2D Hindmarsh--Rose-type system is given by:
\begin{equation}
\begin{aligned}
\dot{x} &= y - ax^3 + bx^2 + I, \\
\dot{y} &= c - dx^2 - y,
\end{aligned}
\end{equation}
where \(x\) represents the membrane potential, \(y\) is the recovery variable, and \(I\) is the external input current. The parameters \(a, b, c, d\) control the intrinsic nonlinearities.

To obtain a mathematically tractable model for global phase-portrait analysis, we fix the parameters as follows:
\begin{itemize}
    \item \(a = 1\)
    \item \(c = 1\)
    \item \(b = d = k\) (where \(k\) is a fixed positive constant, biologically motivated by the standard Hindmarsh--Rose baseline \(k=5\))
\end{itemize}

The system reduces to:
\begin{equation}
\begin{aligned}
\dot{x} &= y - x^3 + kx^2 + I, \\
\dot{y} &= 1 - kx^2 - y.
\end{aligned} \tag{S}
\label{eq:simplified}
\end{equation}
Here, \(I\) is fixed and $k$ treated as the bifurcation parameter.

\section{Hopf Bifurcation and Preliminary Analysis of Limit Cycles}

We consider the following planar system
\begin{equation}
\label{eq:model}
\begin{cases}
\dot{x}=y-x^3+kx^2+I,\\[1mm]
\dot{y}=1-kx^2-y,
\end{cases}
\end{equation}
where $k$ and $I$ are real parameters. The aim of this section is to
investigate the occurrence of Hopf bifurcation and derive a suitable
formulation for the study of periodic solutions.

\subsection{The equilibrium point and Hopf bifurcation}

The equilibrium points of system \eqref{eq:model} satisfy

\[
y-x^3+kx^2+I=0,
\]

and

\[
1-kx^2-y=0.
\]

From the second equation we obtain

\[
y=1-kx^2.
\]

Substituting this expression into the first equation gives

\[
1-kx^2-x^3+kx^2+I=0,
\]

and consequently

\[
x^3=1+I.
\]

Hence system \eqref{eq:model} has a unique equilibrium point

\begin{equation}
\label{eq:equilibrium}
E_*=(x_*,y_*),
\end{equation}

where

\begin{equation}
\label{eq:xstar}
x_*=(1+I)^{1/3},
\end{equation}

and

\begin{equation}
\label{eq:ystar}
y_*=1-kx_*^2 .
\end{equation}

The Jacobian matrix of system \eqref{eq:model} is

\[
J(x,y)=
\begin{pmatrix}
-3x^2+2kx&1\\
-2kx&-1
\end{pmatrix}.
\]

At the equilibrium point $E_*$ we have

\[
J(E_*)=
\begin{pmatrix}
-3x_*^2+2kx_*&1\\
-2kx_*&-1
\end{pmatrix}.
\]

The trace and determinant of this matrix are

\begin{equation}
\label{eq:trace}
T=-3x_*^2+2kx_*-1,
\end{equation}

and

\begin{equation}
\label{eq:det}
D=
(-3x_*^2+2kx_*)(-1)+2kx_*.
\end{equation}

Therefore,

\[
D=3x_*^2.
\]

The necessary conditions for a Hopf bifurcation are

\[
T=0,\qquad D>0.
\]

The trace condition gives

\[
-3x_*^2+2kx_*-1=0,
\]

and therefore the Hopf curve is

\begin{equation}
\label{eq:Hopfcurve}
k=k_H=
\frac{3x_*^2+1}{2x_*}.
\end{equation}

Since

\[
D=3x_*^2,
\]

we have

\[
D>0
\]

provided that

\[
x_*\neq0,
\]

or equivalently,

\[
I\neq -1.
\]

At the Hopf point, the eigenvalues of $J(E_*)$ are

\[
\lambda_{1,2}=\pm i\omega,
\]

where

\begin{equation}
\label{eq:omega}
\omega=\sqrt{3}|x_*|.
\end{equation}


It remains to verify the transversality condition. Since

\[
T(k)=-3x_*^2+2kx_*-1,
\]

we obtain

\[
\frac{dT}{dk}=2x_*.
\]

Because

\[
\operatorname{Re}(\lambda)=\frac{T}{2},
\]

we have

\[
\frac{d}{dk}\operatorname{Re}(\lambda)
=
x_*.
\]

Therefore,

\[
\frac{d}{dk}\operatorname{Re}(\lambda)
\neq0
\]

whenever $I\neq-1$. Consequently, system \eqref{eq:model}
undergoes a Hopf bifurcation on the curve \eqref{eq:Hopfcurve}.


\subsection{The first Lyapunov coefficient}

To determine the type of Hopf bifurcation, we translate the equilibrium
point $E_*$ to the origin by introducing

\[
x=x_*+u,\qquad y=y_*+v .
\]

Then system \eqref{eq:model} becomes

\begin{equation}
\label{eq:shifted}
\begin{cases}
\dot{u}=(-3x_*^2+2kx_*)u+v
 +(k-3x_*)u^2-u^3,\\[1mm]
\dot{v}=-2kx_*u-v-ku^2 .
\end{cases}
\end{equation}

At the Hopf point $k=k_H$, using

\[
2k_Hx_*=3x_*^2+1,
\]

the linear part of \eqref{eq:shifted} is

\[
A=
\begin{pmatrix}
1&1\\
-(3x_*^2+1)&-1
\end{pmatrix}.
\]

The nonlinear terms are

\[
F_2(u,v)=
\begin{pmatrix}
(k_H-3x_*)u^2\\
-k_Hu^2
\end{pmatrix},
\]

and

\[
F_3(u,v)=
\begin{pmatrix}
-u^3\\
0
\end{pmatrix}.
\]

Using the standard formula for the first Lyapunov coefficient of a
planar Hopf bifurcation, after substituting the eigenvectors of $A$ and
simplifying the resulting expression, we obtain

\begin{equation}
\label{eq:l1}
l_1=
-\frac{3x_*^2+1}
{8\sqrt{3}\,x_*^2}.
\end{equation}

Since

\[
3x_*^2+1>0,
\]

we conclude that

\[
l_1<0 .
\]

Hence the Hopf bifurcation is supercritical and a stable limit cycle
bifurcates from the equilibrium point when the parameter crosses the
Hopf curve.

So, until now we have: 

{theorem:}
\label{thm:hopf_limit_cycle}
Consider the planar system
\begin{equation}
\label{eq:system}
\begin{cases}
\dot{x}=y-x^3+kx^2+I,\\[1mm]
\dot{y}=1-kx^2-y,
\end{cases}
\end{equation}
where $k$ and $I$ are real parameters. Let

\[
x_*=(1+I)^{1/3},
\]

and assume that

\[
I\neq -1.
\]

Then system \eqref{eq:system} has a unique equilibrium point

\[
E_*=
\left(
x_*,
1-kx_*^2
\right).
\]

Moreover, the equilibrium undergoes a Hopf bifurcation at the
parameter value

\begin{equation}
\label{eq:hopf}
k=k_H=
\frac{3x_*^2+1}{2x_*}.
\end{equation}

Indeed, at $k=k_H$ the Jacobian matrix of system \eqref{eq:system}
at $E_*$ has a pair of purely imaginary eigenvalues

\[
\lambda_{1,2}=\pm i\omega ,
\]

where

\[
\omega=\sqrt{3}|x_*|,
\]

and the transversality condition

\[
\frac{d}{dk}\operatorname{Re}(\lambda(k))
\bigg|_{k=k_H}\neq0
\]

is satisfied.

Furthermore, after translating the equilibrium point to the origin,
the first Lyapunov coefficient at the Hopf point is

\begin{equation}
\label{eq:l1}
l_1=
-\frac{3x_*^2+1}
{8\sqrt{3}\,x_*^2}.
\end{equation}

Consequently,

\[
l_1<0,
\]

and therefore the Hopf bifurcation is supercritical. Hence, for
parameter values sufficiently close to and beyond the Hopf curve,
system \eqref{eq:system} possesses a stable limit cycle bifurcating
from the equilibrium point $E_*$.
\end{theorem}

{proof:}
The equilibrium point is obtained by solving
\[
y-x^3+kx^2+I=0,
\qquad
1-kx^2-y=0,
\]
which gives
\[
x_*^3=1+I,
\qquad
y_*=1-kx_*^2.
\]

The Jacobian matrix is

\[
J(x,y)=
\begin{pmatrix}
-3x^2+2kx&1\\
-2kx&-1
\end{pmatrix}.
\]

Therefore, at $E_*$,

\[
\operatorname{tr}(J)
=
-3x_*^2+2kx_*-1,
\]

and

\[
\det(J)=3x_*^2.
\]

The Hopf condition
\[
\operatorname{tr}(J)=0
\]
gives

\[
k_H=\frac{3x_*^2+1}{2x_*}.
\]

Since $I\neq-1$, we have $x_*\neq0$ and hence

\[
\det(J)=3x_*^2>0.
\]

Thus the eigenvalues are purely imaginary at the Hopf point.

Furthermore,

\[
\frac{d}{dk}\operatorname{Re}(\lambda)
=
\frac12
\frac{d}{dk}\operatorname{tr}(J)
=
x_* \neq0,
\]

which verifies the transversality condition.

Finally, applying the standard formula for the first Lyapunov
coefficient of a planar Hopf bifurcation to the translated system gives

\[
l_1=
-\frac{3x_*^2+1}{8\sqrt3\,x_*^2}.
\]

Since $x_*\neq0$, we have $l_1<0$. Hence the Hopf bifurcation is
supercritical and a stable periodic orbit is generated from the
equilibrium point.
\end{proof}


\subsection{Transformation into a Lienard system}

For the global study of periodic solutions, we rewrite system
\eqref{eq:model} as a second-order differential equation.

From the second equation,

\[
\dot y=1-kx^2-y,
\]

we obtain

\[
y=1-kx^2-\dot y .
\]

Substitution into the first equation gives

\[
\dot{x}=1+I-x^3-\dot y .
\]

Differentiating this relation yields

\[
\ddot{x}=-3x^2\dot{x}-\ddot{y}.
\]

On the other hand,

\[
\ddot{y}+2kx\dot{x}+\dot y=0.
\]

Combining the above relations gives the following second-order equation:

\begin{equation}
\label{eq:lienard}
\ddot{x}
+
(3x^2-2kx+1)\dot{x}
+
x^3-(1+I)=0 .
\end{equation}

Therefore, system \eqref{eq:model} is equivalent to the Lienard system

\begin{equation}
\label{eq:lienardplanar}
\begin{cases}
\dot{x}=z,\\[1mm]
\dot{z}=-(3x^2-2kx+1)z-x^3+(1+I).
\end{cases}
\end{equation}

The damping function of the Lienard system is

\[
f(x)=3x^2-2kx+1,
\]

and its primitive function is

\[
F(x)=\int_0^x f(s)\,ds
=x^3-kx^2+x .
\]

The damping function changes sign if and only if

\[
\Delta_f=(-2k)^2-12>0,
\]

which gives

\begin{equation}
\label{eq:dampingcondition}
k^2>3 .
\end{equation}

Therefore, the parameter region satisfying \eqref{eq:dampingcondition}
is the natural region where self-sustained oscillations and limit cycles
may occur.

The above transformation provides the basis for the global analysis of
periodic solutions.





\newpage

pppppppppppppppp


\subsection{Trace and Determinant}
The trace \(T\) and determinant \(D\) of \(J^*\) are:
\begin{align}
T &= \operatorname{tr}(J^*) = -3u^2 + 2ku - 1, \\
D &= \det(J^*) = (-3u^2 + 2ku)(-1) - (1)(-2ku) = 3u^2 - 2ku + 2ku = 3u^2.
\end{align}
Notice that \(D = 3u^2 > 0\) for all \(I > -1\). Therefore, the equilibrium is \textbf{never a saddle}. It is either a node or a focus, and its stability is determined entirely by the sign of the trace \(T\).

\subsection{Classification of the Singular Point}
\begin{itemize}
    \item If \(T^2 - 4D > 0\): Node (stable if \(T < 0\), unstable if \(T > 0\)).
    \item If \(T^2 - 4D < 0\): Focus (stable if \(T < 0\), unstable if \(T > 0\)).
    \item If \(T = 0\) and \(D > 0\): Center (linearization) or weak focus.
\end{itemize}

\section{Poincar\'e Compactification and the Dynamics at Infinity}

\label{sec:infinity}

The vector field associated with system \eqref{eq:model} is polynomial of
degree three. Therefore, the behavior of solutions at infinity can be
studied by means of the Poincar\'e compactification. We analyze the
compactified vector field in two local charts corresponding to the
directions of the coordinate axes.


\subsection{Local chart \(U_1\)}

In the first chart, we introduce the coordinates

\[
x=\frac{1}{u},\qquad y=\frac{v}{u}.
\]

Since the degree of the vector field is three, the time transformation

\[
d\tau=\frac{dt}{u^{2}}
\]

gives the compactified system

\begin{equation}
\label{eq:U1}
\begin{aligned}
\dot u&=-u^3v+u-ku^2-Iu^4,\\
\dot v&=u^3-ku-u^2v-u^2v^2+v-kuv-Iu^3v .
\end{aligned}
\end{equation}

The line at infinity in this chart is

\[
u=0 .
\]

Restricting \eqref{eq:U1} to this invariant line gives

\[
\dot u=0,\qquad
\dot v=v .
\]

Hence the only singular point on the equator in \(U_1\) is

\[
P_1=(0,0).
\]

The Jacobian matrix at \(P_1\) is

\[
J_{U_1}(P_1)
=
\begin{pmatrix}
1&0\\
-k&1
\end{pmatrix}.
\]

Both eigenvalues are equal to one. Consequently,

\[
\boxed{
P_1\ \text{is a hyperbolic unstable node at infinity}.
}
\]

\subsection{Local chart \(U_2\)}

The second chart is defined by

\[
x=\frac{u}{v},\qquad y=\frac1v .
\]

After the rescaling

\[
d\tau=\frac{dt}{v^2},
\]

we obtain

\begin{equation}
\label{eq:U2}
\begin{aligned}
\dot u={}&
-u^3+ku^3v+ku^2v+uv^2-uv^3+v^2+Iv^3,
\\
\dot v={}&
ku^2v^2+v^3-v^4 .
\end{aligned}
\end{equation}

The infinity line is

\[
v=0 .
\]

On this line,

\[
\dot u=-u^3,\qquad \dot v=0 .
\]

Thus the unique singular point in this chart is

\[
P_2=(0,0).
\]

The Jacobian matrix satisfies

\[
J_{U_2}(P_2)=
\begin{pmatrix}
0&0\\
0&0
\end{pmatrix}.
\]

Therefore \(P_2\) is a linearly zero singular point. Since the linear
part gives no information about the local phase portrait, a
desingularization by a blow-up technique is required.

%=========================================================
\subsection{Blow-up analysis of the degenerate point in $U_2$}
\label{subsec:blowup-U2}
%=========================================================

In the local chart $U_2$, the compactified vector field is
\begin{equation}
\label{eq:U2-system}
\begin{aligned}
\dot u
&=
-u^3
+ku^3v
+ku^2v
+uv^2
-uv^3
+v^2
+Iv^3,
\\
\dot v
&=
ku^2v^2+v^3-v^4.
\end{aligned}
\end{equation}
The line at infinity is $v=0$, and the restriction of
\eqref{eq:U2-system} to this line is
\[
\dot u=-u^3,
\qquad
\dot v=0.
\]
Hence the only singular point on the equator in this chart is
\[
P_2=(0,0).
\]
Moreover,
\[
D X_{U_2}(P_2)
=
\begin{pmatrix}
0&0\\
0&0
\end{pmatrix},
\]
so that $P_2$ is completely degenerate and linearization gives no
information about its local phase portrait.

We therefore resolve this singularity by means of directional blow-ups.

%---------------------------------------------------------
\subsubsection{First directional blow-up}
%---------------------------------------------------------

Let
\[
P(u,v)
=
-u^3+ku^3v+ku^2v+uv^2-uv^3+v^2+Iv^3,
\]
and
\[
Q(u,v)
=
ku^2v^2+v^3-v^4.
\]
The lowest-order nonzero homogeneous terms are
\[
P_2(u,v)=v^2,
\qquad
Q_2(u,v)=0.
\]
Consequently, the characteristic polynomial is
\[
\mathcal F(u,v)
=
uQ_2(u,v)-vP_2(u,v)
=
-v^3.
\]
Thus the only characteristic direction is
\[
v=0.
\]
It is therefore natural to perform the directional blow-up
\begin{equation}
\label{eq:first-blowup}
v=uz,
\end{equation}
where $z=v/u$ represents the slope with respect to the $u$-axis.

Since
\[
\dot v=z\dot u+u\dot z,
\]
substitution of $v=uz$ into the first equation of
\eqref{eq:U2-system} gives
\begin{align*}
\dot u
={}&
-u^3
+ku^4z
+ku^3z
+u^3z^2
-u^4z^3
+u^2z^2
+Iu^3z^3
\\
={}&
u^2
\left(
z^2-u+kuz+uz^2+Iuz^3
+ku^2z-u^2z^3
\right).
\end{align*}
Similarly,
\[
Q(u,uz)
=
ku^4z^2+u^3z^3-u^4z^4.
\]
Therefore,
\[
u\dot z
=
Q(u,uz)-zP(u,uz),
\]
and hence
\[
\dot z
=
uz
\left(
u-z^2-kuz-Iuz^3
\right).
\]
Thus, before desingularization, the blown-up system reads
\[
\begin{aligned}
\dot u
&=
u^2
\left(
z^2-u+kuz+uz^2+Iuz^3
+ku^2z-u^2z^3
\right),
\\
\dot z
&=
uz
\left(
u-z^2-kuz-Iuz^3
\right).
\end{aligned}
\]
Both equations contain the common factor $u$. Cancelling this factor,
which corresponds to a time rescaling away from the exceptional
divisor, yields the desingularized system
\begin{equation}
\label{eq:first-desingularized}
\boxed{
\begin{aligned}
u'
&=
u
\left(
z^2-u+kuz+uz^2+Iuz^3
+ku^2z-u^2z^3
\right),
\\
z'
&=
z
\left(
u-z^2-kuz-Iuz^3
\right).
\end{aligned}}
\end{equation}
Here the prime denotes differentiation with respect to the
desingularized time.

The exceptional divisor is
\[
u=0.
\]
Restricting \eqref{eq:first-desingularized} to $u=0$ gives
\[
u'=0,
\qquad
z'=-z^3.
\]
Hence the only singular point on the exceptional divisor is again
\[
(u,z)=(0,0).
\]
Moreover,
\[
D X_1(0,0)
=
\begin{pmatrix}
0&0\\
0&0
\end{pmatrix}.
\]
Thus the first directional blow-up does not completely resolve the
degeneracy.

The lowest-order homogeneous terms of
\eqref{eq:first-desingularized} are
\[
P^{(1)}_2(u,z)=-u^2,
\qquad
Q^{(1)}_2(u,z)=uz.
\]
Therefore the corresponding characteristic polynomial is
\[
\mathcal F_1(u,z)
=
uQ^{(1)}_2(u,z)-zP^{(1)}_2(u,z)
=
2u^2z.
\]
Hence the new characteristic directions are
\[
z=0
\qquad\text{and}\qquad
u=0.
\]
We consequently study these two directions in complementary blow-up
charts.

%---------------------------------------------------------
\subsubsection{Second blow-up: the direction $z=0$}
%---------------------------------------------------------

To resolve the direction $z=0$, introduce
\begin{equation}
\label{eq:second-blowup-chart1}
z=uw.
\end{equation}
Since
\[
\dot z=w\dot u+u\dot w,
\]
substitution into \eqref{eq:first-desingularized} gives
\[
u'
=
u^2
\left[
-1
+kuw
+uw^2
+ku^2w
+u^2w^2
+Iu^3w^3
-u^4w^3
\right],
\]
while
\[
w'
=
uw
\left[
2
-2kuw
-2uw^2
-ku^2w
-u^2w^2
-2Iu^3w^3
+u^4w^3
\right].
\]
After cancelling the common factor $u$, the desingularized system is
\begin{equation}
\label{eq:second-desing-chart1}
\boxed{
\begin{aligned}
u'
&=
u
\left[
-1
+kuw
+uw^2
+ku^2w
+u^2w^2
+Iu^3w^3
-u^4w^3
\right],
\\
w'
&=
w
\left[
2
-2kuw
-2uw^2
-ku^2w
-u^2w^2
-2Iu^3w^3
+u^4w^3
\right].
\end{aligned}}
\end{equation}
On the exceptional divisor $u=0$,
\[
u'=0,
\qquad
w'=2w.
\]
Thus the only singular point is $(0,0)$.

The Jacobian matrix at this point is
\[
D X_2(0,0)
=
\begin{pmatrix}
-1&0\\
0&2
\end{pmatrix}.
\]
Hence
\[
\lambda_1=-1,
\qquad
\lambda_2=2.
\]
Consequently,
\[
\boxed{
(0,0)\text{ is a hyperbolic saddle in this blow-up chart}.
}
\]
Thus the characteristic direction $z=0$ is completely resolved.

%---------------------------------------------------------
\subsubsection{Second blow-up: the complementary direction $u=0$}
%---------------------------------------------------------

The direction $u=0$ is not represented by the previous chart.
We therefore introduce the complementary directional blow-up
\begin{equation}
\label{eq:second-blowup-chart2}
u=zr.
\end{equation}
Since
\[
\dot u=r\dot z+z\dot r,
\]
substitution into \eqref{eq:first-desingularized} gives, after
cancelling a common factor $z$,
\begin{equation}
\label{eq:second-desing-chart2}
\boxed{
\begin{aligned}
z'
&=
z\left(
r-z-kzr-Iz^3r
\right),
\\
r'
&=
r\left(
2Iz^3r
+kz^2r^2
+2kzr
-z^4r^2
+z^2r
+2z
-2r
\right).
\end{aligned}}
\end{equation}
The exceptional divisor is now $z=0$. Restriction to this divisor
gives
\[
z'=0,
\qquad
r'=-2r^2.
\]
Thus the unique singularity on the exceptional divisor is again
\[
(z,r)=(0,0),
\]
which remains degenerate.

The quadratic homogeneous part of
\eqref{eq:second-desing-chart2} is
\[
P^{(2)}_2(z,r)=z(r-z),
\]
and
\[
Q^{(2)}_2(z,r)=2r(z-r).
\]
Hence the characteristic polynomial is
\begin{align*}
\mathcal F_2(z,r)
&=
zQ^{(2)}_2(z,r)-rP^{(2)}_2(z,r)
\\
&=
2zr(z-r)-rz(r-z)
\\
&=
3zr(z-r).
\end{align*}
Therefore the characteristic directions are
\[
z=0,
\qquad
r=0,
\qquad
r=z.
\]

%---------------------------------------------------------
\subsubsection{Third blow-up along the direction $r=z$}
%---------------------------------------------------------

To resolve the nontrivial characteristic direction $r=z$, set
\begin{equation}
\label{eq:third-blowup}
r=zs.
\end{equation}
Since
\[
\dot r=s\dot z+z\dot s,
\]
substitution of \eqref{eq:third-blowup} into
\eqref{eq:second-desing-chart2} and cancellation of the common factor
$z$ yield
\begin{equation}
\label{eq:third-desing}
\boxed{
\begin{aligned}
z'
&=
z
\left(
s-1-ksz-Isz^3
\right),
\\
s'
&=
s
\left(
3-3s
+3ksz
+sz^2
+3Isz^3
+ks^2z^3
-s^2z^5
\right).
\end{aligned}}
\end{equation}

The exceptional divisor is $z=0$, and on it
\[
z'=0,
\qquad
s'=3s(1-s).
\]
Consequently, two singular points appear:
\[
R_0=(0,0),
\qquad
R_1=(0,1).
\]

At $R_0=(0,0)$,
\[
D X_3(R_0)
=
\begin{pmatrix}
-1&0\\
0&3
\end{pmatrix},
\]
so that
\[
\lambda_1=-1,
\qquad
\lambda_2=3.
\]
Therefore,
\[
\boxed{R_0\text{ is a hyperbolic saddle}.}
\]

At $R_1=(0,1)$, the Jacobian has eigenvalues
\[
\boxed{0,\,-3}.
\]
Hence $R_1$ is semi-hyperbolic and its local behavior must be
determined from the dynamics on its center manifold.

%---------------------------------------------------------
%=========================================================
%---------------------------------------------------------
\subsubsection{Complementary chart of the third blow-up}
\label{subsubsec:third-complementary}
%---------------------------------------------------------

The preceding chart
\[
r=zs
\]
resolves the characteristic directions $r=0$ and $r=z$, but it does
not represent the remaining direction $z=0$. To complete the third
blow-up, we therefore introduce the complementary chart
\begin{equation}
\label{eq:third-complementary}
z=rq.
\end{equation}

Recall that, before this third blow-up, the desingularized vector field
is
\begin{equation}
\label{eq:second-desing-chart2}
\begin{aligned}
z'
&=
z\left(
r-z-kzr-Iz^3r
\right),
\\
r'
&=
r\left(
2Iz^3r
+kz^2r^2
+2kzr
-z^4r^2
+z^2r
+2z
-2r
\right).
\end{aligned}
\end{equation}

Since
\[
z=rq,
\]
we have
\[
z'=q r'+r q'.
\]
Substituting $z=rq$ into \eqref{eq:second-desing-chart2} and cancelling
the common factor $r$ gives
\begin{equation}
\label{eq:third-complementary-desing}
\begin{aligned}
q'
&=
q\Big(
3-3q
-3krq
-r^2q^2
-kr^3q^2
-3Ir^3q^3
+r^5q^4
\Big),
\\
r'
&=
r\Big(
-2+2q
+2krq
+r^2q^2
+kr^3q^2
+2Ir^3q^3
-r^5q^4
\Big).
\end{aligned}
\end{equation}

The exceptional divisor in this chart is
\[
r=0.
\]
Restricting \eqref{eq:third-complementary-desing} to $r=0$ yields
\[
q'=3q(1-q),
\qquad
r'=0.
\]
Hence the singular points on the exceptional divisor are
\[
Q_0=(0,0),
\qquad
Q_1=(1,0).
\]

At $Q_0$,
\[
D X(Q_0)
=
\begin{pmatrix}
3&0\\
0&-2
\end{pmatrix},
\]
and therefore
\[
\lambda_1=3,
\qquad
\lambda_2=-2.
\]
Thus
\[
\boxed{Q_0\text{ is a hyperbolic saddle}.}
\]

The point $Q_1$ does not represent a new singular direction. Indeed,
on the overlap of the two third blow-up charts
\[
r=zs,
\qquad
z=rq,
\]
we obtain
\[
sq=1.
\]
Therefore
\[
q=1
\quad\Longleftrightarrow\quad
s=1,
\]
so $Q_1=(1,0)$ corresponds precisely to the semi-hyperbolic point
\[
R_1=(0,1)
\]
already obtained in the chart $r=zs$.

Consequently, the two third-blow-up charts together cover all three
characteristic directions
\[
z=0,\qquad r=0,\qquad r=z.
\]
The direction $z=0$ gives the hyperbolic saddle $Q_0$, while the
directions represented in the first chart give the hyperbolic saddle
$R_0$ and the semi-hyperbolic point $R_1$. The latter is analyzed
below by a center-manifold reduction.

%=========================================================
\subsubsection{Center-manifold reduction at $R_1=(0,1)$}
\label{subsubsec:center-R1}
%=========================================================

We now analyze the remaining nonhyperbolic point
\[
R_1=(0,1)
\]
of \eqref{eq:third-desing}. Introduce
\[
s=1+h,
\]
so that $R_1$ is translated to $(z,h)=(0,0)$. System
\eqref{eq:third-desing} becomes
\begin{equation}
\label{eq:R1-translated}
\begin{aligned}
z'
&=
z\Big[
h-k(1+h)z-I(1+h)z^3
\Big],
\\
h'
&=
(1+h)\Big[
-3h
+3k(1+h)z
+(1+h)z^2
\\
&\hspace{30mm}
+3I(1+h)z^3
+k(1+h)^2z^3
-(1+h)^2z^5
\Big].
\end{aligned}
\end{equation}

The Jacobian matrix at the origin is
\[
D X(0,0)
=
\begin{pmatrix}
0&0\\
3k&-3
\end{pmatrix}.
\]
Its eigenvalues are
\[
\lambda_1=0,
\qquad
\lambda_2=-3.
\]
Thus $R_1$ is semi-hyperbolic.

For the zero eigenvalue, an eigenvector $(a,b)^T$ satisfies
\[
3ka-3b=0,
\]
and hence
\[
b=ka.
\]
Therefore the center eigenspace is
\[
E^c=\{(z,h):h=kz\}.
\]
Consequently, a local center manifold may be represented as
\[
h=H(z),
\]
where
\begin{equation}
\label{eq:center-tangency}
H(0)=0,
\qquad
H'(0)=k.
\end{equation}
Notice that $H'(0)$ is not zero in the present coordinates.

We seek an expansion
\begin{equation}
\label{eq:center-expansion-assumption}
H(z)
=
a_1z+a_2z^2+a_3z^3+O(z^4).
\end{equation}
The invariance equation is
\begin{equation}
\label{eq:center-invariance}
H'(z)\,z'
=
h'
\qquad\text{on }h=H(z).
\end{equation}
Substitution of \eqref{eq:center-expansion-assumption} into
\eqref{eq:center-invariance} and comparison of like powers of $z$
gives, at first order,
\[
3a_1-3k=0,
\]
and hence
\[
a_1=k.
\]
At second order,
\[
4a_1^2-7ka_1+3a_2-1=0.
\]
Using $a_1=k$, this becomes
\[
3a_2-3k^2-1=0,
\]
so
\[
a_2=k^2+\frac13.
\]
At third order,
\[
-3I
-4a_1^2k
+9a_1a_2
-2a_1
-8ka_2
+3a_3-k
=0.
\]
Substituting the preceding values of $a_1$ and $a_2$ yields
\[
3a_3
-3I
-3k^3
-\frac{8k}{3}
=0.
\]
Therefore,
\[
a_3
=
I+k^3+\frac{8k}{9}.
\]

Hence the center manifold through $R_1$ is
\begin{equation}
\label{eq:center-manifold-R1}
\boxed{
H(z)
=
kz
+
\left(k^2+\frac13\right)z^2
+
\left(
I+k^3+\frac{8k}{9}
\right)z^3
+
O(z^4).
}
\end{equation}
Equivalently, in the $(z,s)$ coordinates,
\[
s
=
1+kz
+
\left(k^2+\frac13\right)z^2
+
\left(
I+k^3+\frac{8k}{9}
\right)z^3
+
O(z^4).
\]

Substituting \eqref{eq:center-manifold-R1} into the first equation of
\eqref{eq:R1-translated} gives the reduced equation on the center
manifold:
\begin{equation}
\label{eq:reduced-R1}
\boxed{
z'
=
\frac13z^3
+
\frac{5k}{9}z^4
+
O(z^5).
}
\end{equation}
In particular, the first nonvanishing term is
\[
\frac13z^3.
\]
Therefore, for $|z|$ sufficiently small,
\[
z'>0
\quad\text{if }z>0,
\qquad
z'<0
\quad\text{if }z<0.
\]
Thus the flow on the center manifold is repelling from $R_1$ on both
sides. On the other hand, the transverse eigenvalue is
\[
-3<0,
\]
and therefore the transverse direction is attracting.

Consequently, $R_1$ has saddle-type local behavior for the
desingularized vector field: the center manifold provides the
unstable direction, while the transverse direction is stable.

%=========================================================
\subsubsection{Geometric interpretation of the repeated blow-ups}
\label{subsubsec:blowup-geometry}
%=========================================================

The transformations leading to the chart containing $R_1$ are
\[
v=uz,
\qquad
u=zr,
\qquad
r=zs.
\]
Combining them gives the exact blow-down relations
\begin{equation}
\label{eq:combined-blowdown}
\boxed{
u=z^2s,
\qquad
v=z^3s.
}
\end{equation}
These relations hold throughout the third blow-up chart.

For trajectories approaching $R_1=(0,1)$ we have
\[
s\longrightarrow1.
\]
Hence, specifically near $R_1$,
\[
u\sim z^2,
\qquad
v\sim z^3.
\]
It follows that
\[
\frac{v^2}{u^3}
=
\frac{z^6s^2}{z^6s^3}
=
\frac1s
\longrightarrow1,
\]
and consequently
\begin{equation}
\label{eq:cusp-balance}
\boxed{
v^2\sim u^3
\qquad
\text{along trajectories approaching }R_1.
}
\end{equation}

This weighted balance is already suggested by the leading terms of
the original $U_2$ equation,
\[
\dot u
=
v^2-u^3+\text{higher-order terms}.
\]
The monomials $v^2$ and $u^3$ have equal weighted degree if
\[
\operatorname{wt}(u)=2,
\qquad
\operatorname{wt}(v)=3.
\]
This motivates a quasi-homogeneous blow-up of weight $(2,3)$.


\subsubsection{Alternative quasi-homogeneous $(2,3)$-blow-up}
\label{subsubsec:weighted-blowup}
%=========================================================

The repeated ordinary blow-ups above already provide a complete local
desingularization. Nevertheless, the asymptotic balance
\[
v^2\sim u^3
\]
shows that the singularity has a natural quasi-homogeneous structure.
Following the standard quasi-homogeneous desingularization procedure
\cite{DumortierLlibreArtes2006,BrunellaMiari1990,Pelletier1995},Doi: 10.1016/0022-0396(90)90120-E
we assign the weights
\[
\operatorname{wt}(u)=2,
\qquad
\operatorname{wt}(v)=3.
\]
Indeed, the dominant monomials $u^3$ and $v^2$ then have the same
weighted degree,
\[
3\,\operatorname{wt}(u)
=
2\,\operatorname{wt}(v)
=
6.
\]
Accordingly, in the positive-$u$ weighted chart we introduce
\begin{equation}
\label{eq:weighted-blowup}
u=r^2,
\qquad
v=r^3w,
\qquad r\geq0.
\end{equation}
Then
\[
u^3=r^6,
\qquad
v^2=r^6,
\]
so the two dominant monomials are balanced at the same order.

Differentiating gives
\[
\dot u=2r\dot r,
\qquad
\dot v=3r^2w\dot r+r^3\dot w.
\]
Substitution into the original $U_2$ vector field yields
\[
\dot r
=
\frac{r^5}{2}
\left[
w^2-1
+krw
+r^2w^2
+kr^3w
+Ir^3w^3
-r^5w^3
\right]
\]
and
\[
\dot w
=
\frac{r^4w}{2}
\left[
3-3w^2
-3krw
-r^2w^2
-kr^3w
-3Ir^3w^3
+r^5w^3
\right].
\]
After division by the common factor $r^4$, the desingularized system
is
\begin{equation}
\label{eq:weighted-desing}
\boxed{
\begin{aligned}
r'
&=
\frac{r}{2}
\left[
w^2-1
+krw
+r^2w^2
+kr^3w
+Ir^3w^3
-r^5w^3
\right],
\\
w'
&=
\frac{w}{2}
\left[
3-3w^2
-3krw
-r^2w^2
-kr^3w
-3Ir^3w^3
+r^5w^3
\right].
\end{aligned}}
\end{equation}

On the exceptional divisor $r=0$,
\[
r'=0,
\qquad
w'=\frac32w(1-w^2).
\]
Hence
\[
W_0=(0,0),
\qquad
W_+=(0,1),
\qquad
W_-=(0,-1)
\]
are the singular points in this chart.

At $W_0$,
\[
D X_W(W_0)
=
\begin{pmatrix}
-\frac12&0\\
0&\frac32
\end{pmatrix},
\]
and therefore $W_0$ is a hyperbolic saddle.

At both $W_+$ and $W_-$ the eigenvalues are
\[
0,\qquad -3.
\]
Thus these two points are semi-hyperbolic. Their center-manifold
reductions may be obtained explicitly. Writing
\[
w=\varepsilon+\eta,
\qquad
\varepsilon\in\{1,-1\},
\]
one finds
\[
\eta
=
-\frac{k}{2}r
+
\varepsilon
\left(
\frac{k^2}{8}-\frac16
\right)r^2
+
\left(
-\frac I2+\frac{k}{18}
\right)r^3
+
O(r^4),
\]
and the reduced radial equation is
\begin{equation}
\label{eq:weighted-reduced}
\boxed{
r'
=
\frac13r^3
+
\varepsilon\frac{k}{18}r^4
+
O(r^5).
}
\end{equation}
Thus the leading center dynamics is again repelling from the
exceptional divisor, whereas the transverse eigenvalue $-3$ is
attracting.

Moreover, since
\[
u=r^2,
\qquad
v=r^3w,
\]
we have
\[
\frac{v}{u^{3/2}}=w.
\]
Consequently,
\[
W_+:\quad v\sim u^{3/2},
\qquad
W_-:\quad v\sim -u^{3/2}.
\]
The weighted blow-up therefore separates directly the two real
branches associated with the balance
\[
v^2\sim u^3.
\]

The chart \eqref{eq:weighted-blowup} covers only the sector $u\geq0$.
For completeness, the negative-$u$ weighted chart is
\[
u=-r^2,
\qquad
v=r^3w,
\qquad r\geq0.
\]
After the same desingularization, one obtains
\begin{equation}
\label{eq:weighted-negative-u}
\begin{aligned}
r'
&=
\frac r2
\left[
-1-w^2
-krw
+r^2w^2
+kr^3w
-Ir^3w^3
-r^5w^3
\right],
\\
w'
&=
\frac w2
\left[
3+3w^2
+3krw
-r^2w^2
-kr^3w
+3Ir^3w^3
+r^5w^3
\right].
\end{aligned}
\end{equation}
On $r=0$,
\[
r'=0,
\qquad
w'=\frac32w(1+w^2),
\]
so the only real singular point is
\[
W_0^-=(0,0).
\]
Its Jacobian is
\[
D X(W_0^-)
=
\begin{pmatrix}
-\frac12&0\\
0&\frac32
\end{pmatrix},
\]
and hence $W_0^-$ is again a hyperbolic saddle.

The remaining $v$-dominant weighted charts do not generate additional
singular directions. Indeed, using
\[
u=\rho^2\xi,
\qquad
v=\pm\rho^3,
\]
the restriction of the corresponding desingularized vector field to
$\rho=0$ gives
\[
\xi'=1-\xi^3.
\]
Thus the only real singular value is
\[
\xi=1,
\]
which corresponds, on the overlaps, precisely to $W_+$ or $W_-$.
Hence no additional singular points arise from the missing weighted
directions.

%=========================================================
\subsubsection{Conclusion of the blow-up analysis}
\label{subsubsec:blowup-conclusion}
%=========================================================

The point $P_2=(0,0)$ in the compactified chart $U_2$ is completely
degenerate at the linear level. Ordinary directional blow-ups resolve
the characteristic directions into hyperbolic saddles together with
a single semi-hyperbolic point $R_1$. The center-manifold reduction at
$R_1$ gives
\[
z'
=
\frac13z^3
+
O(z^4),
\]
so the center direction is repelling, while the nonzero transverse
eigenvalue is negative. Thus $R_1$ has saddle-type local behavior in
the desingularized vector field.

The exact blow-down relation near this point is
\[
u=z^2s,
\qquad
v=z^3s,
\]
and, since $s\to1$ along trajectories tending to $R_1$,
\[
v^2\sim u^3.
\]
This explains naturally the weight assignment
\[
\operatorname{wt}(u)=2,
\qquad
\operatorname{wt}(v)=3.
\]
The corresponding quasi-homogeneous blow-up provides an alternative
and more direct description of this weighted cusp-type scaling,
separating the two real branches
\[
v\sim\pm u^{3/2}
\]
in the positive-$u$ sector.

Therefore the degeneracy of $P_2$ is associated with a genuine
higher-order directional structure at infinity. Both the repeated
ordinary blow-up and the $(2,3)$-weighted blow-up reveal the same
underlying scaling, while the ordinary blow-up atlas together with
the center-manifold reduction provides a complete local
desingularization of the singular point.









\end{document}